\subsection{SPECIAL POINTS}

Let L be the set of translations belonging to W and let \(\Lambda\) be the set of \(t \in T\) such that the translation \(x \mapsto t + x\) belongs to L. It is immediate that \(\Lambda\) is stable under U(W) and that L is a normal subgroup of W. Since W acts properly on E, the same holds for L, and it follows easily that \(\Lambda\) is a discrete subgroup of T. For any point x of E, denote by \(W_x\) the stabiliser of x in W.

PROPOSITION 9. Let \(a \in E\) . The following conditions are equivalent:

\begin{enumerate}
\def\labelenumi{(\roman{enumi})}
\item
  We have \(\mathrm{W} = \mathrm{W}_{a}.\mathrm{L}\) ;
\item
  The restriction of the homomorphism U to \(W_{a}\) is an isomorphism from \(W_{a}\) to \(\mathrm{U}(\mathrm{W})\) ;
\item
  For every hyperplane \(\mathbf{H} \in \mathfrak{H}\) , there exists a hyperplane \(\mathbf{H}' \in \mathfrak{H}\) parallel to \(\mathbf{H}\) and such that \(a \in \mathbf{H}'\) .
\end{enumerate}

It is clear that (i) \(\Leftrightarrow\) (ii), since L is the kernel of U and \(\mathrm{L} \cap \mathrm{W}_a = \{1\}\) .

Assume (i) and let \(\mathbf{H} \in \mathfrak{H}\) ; then \(s_{\mathbf{H}} \in \mathbf{W}_a\) . L so there exists a vector \(t \in \Lambda\) such that \(a = s_{\mathbf{H}}(a) + t\) ; the vector \(t\) is orthogonal to \(\mathbf{H}\) , and if \(\mathrm{H}' = \mathrm{H} + \frac{1}{2} t\) then \(s_{\mathbf{H}'}(x) = s_{\mathbf{H}}(x) + t\) for all \(x \in \mathbf{E}\) (cf.~§2, no. 4, Prop. 5). Since \(t \in \Lambda\) and \(s_{\mathbf{H}} \in \mathbf{W}\) , we have \(s_{\mathbf{H}'} \in \mathbf{W}\) , and so \(\mathrm{H}' \in \mathfrak{H}\) ; we also have \(a = s_{\mathbf{H}'}(a)\) , and so \(a \in \mathbf{H}'\) . Thus, (i) implies (iii).

Assume (iii). Let \(\mathrm{H} \in \mathfrak{H}\) ; take \(\mathrm{H}'\) as in (iii). Then \(s_{\mathrm{H}'}(a) = a\) , so \(s_{\mathrm{H}'} \in \mathrm{W}_a\) ; since \(\mathrm{H}\) is parallel to \(\mathrm{H}'\) , the element \(w = s_{\mathrm{H}'} s_{\mathrm{H}}\) of \(\mathrm{W}\) is a translation (§2, no. 4, Prop. 5), so \(w \in \mathrm{L}\) ; then \(s_{\mathrm{H}} = s_{\mathrm{H}'} w \in \mathrm{W}_a\) . L. Since \(\mathrm{W}\) is generated by the family \((s_{\mathrm{H}})_{\mathrm{H} \in \mathfrak{H}}\) , it follows that \(\mathrm{W} = \mathrm{W}_a\) . L and hence (iii) implies (i).

DEFINITION 1. A point \(a\) of \(E\) is special for \(W\) if it satisfies the equivalent conditions in Prop. 9.

It is clear that the set of special points of E is stable under W.

PROPOSITION 10. There exists a special point for W.

By Prop. 6 of no. 8, we need only consider the case when W is essential. The group U(W) of automorphisms of T is finite (cf.~no. 6, Th. 3) and \(U(s_{H})\) is an orthogonal reflection for every hyperplane H; moreover, U(W) is generated by the family \((U(s_{H}))_{H \in \mathfrak{H}}\). By Prop. 7 of no. 9, there exists a basis \((e_{i})_{i \in I}\) of T such that the group U(W) is generated by the set of reflections \((s_{i})_{i \in I}\) defined by

\[
s _ {i} (t) = t - 2 (t | e _ {i}). e _ {i}.
\]

The Cor. of Th. 1 of no. 2 shows that every reflection \(s \in \mathrm{U}(\mathrm{W})\) is of the form \(s = \mathrm{U}(s_{\mathrm{H}})\) with \(\mathrm{H} \in \mathfrak{H}\) . We can thus find in \(\mathfrak{H}\) a family of hyperplanes \((\mathbf{H}_i)_{i \in I}\) such that \(s_i = \mathrm{U}(s_{\mathrm{H}_i})\) for all \(i\) . Since the vectors \(e_i\) are linearly independent, there exists \(a \in E\) such that \(a \in H_i\) for all \(i \in I\) . We have \(s_{\mathrm{H}_i} \in W_a\) , so \(\mathrm{U}(\mathrm{W}) = \mathrm{U}(\mathrm{W}_a)\) , which means that \(W = W_a.L\) since \(L\) is the kernel of \(U\) . Thus, \(a\) is special.

When W is finite and essential, there is only one special point for W, namely the unique point invariant under W. Thus, the consideration of special points is interesting mainly when W is infinite.

PROPOSITION 11. Assume that W is essential. Let a be a special point for W. The chambers relative to \(W_{a}\) are the open simplicial cones with vertex a. For every chamber \(C'\) relative to \(W_{a}\) , there exists a unique chamber C relative to W contained in \(C'\) and such that \(a \in \overline{C}\) . The union of the \(w'(\overline{C})\) for \(w' \in W_{a}\) is a closed neighbourhood of a in E. Every wall of \(C'\) is a wall of C. If W is infinite and irreducible, the walls of C are the walls of \(C'\) together with an affine hyperplane not parallel to the walls of \(C'\) .

Let \(\mathfrak{H}'\) be the set of \(H \in \mathfrak{H}\) containing \(a\) . The group \(W_a\) is generated by the \(s_H\) for \(H \in \mathfrak{H}'\) (no. 3, Prop. 2). The chambers relative to \(W_a\) are the open simplicial cones with vertex \(a\) (no. 9, Prop. 7). Let \(C'\) be such a chamber and let \(U\) be a non-empty open ball with centre \(a\) not meeting any element of \(\mathfrak{H} - \mathfrak{H}'\) . Since \(a \in \overline{C'}\) , there exists a point \(b\) in \(U \cap C'\) . Then \(b \notin H\) for all \(H \in \mathfrak{H}\) , so \(b\) belongs to a chamber \(C\) relative to \(\mathfrak{H}\) . Since \(\mathfrak{H}' \subset \mathfrak{H}\) , we have \(C \subset C'\) . The set \(U \cap C'\) does not meet any \(H \in \mathfrak{H}\) and is convex, so \(U \cap C' \subset C\) ; thus \(a \in \overline{C}\) . Conversely, let \(C_1\) be a chamber relative to \(W\) contained in \(C'\) and such that \(a \in \overline{C}_{1}\) ; then \(C_{1}\) meets U and \(U \cap C_{1} \subset U \cap C' = U \cap C\) ; the chambers C and \(C_{1}\) , having a point in common, must coincide. For any \(w' \in W_{a}\) , we have \(w'(U) = U\) , so

\[
\mathrm{U} \cap w ^ {\prime} (\mathrm{C}) = w ^ {\prime} (\mathrm{U} \cap \mathrm{C}) = w ^ {\prime} (\mathrm{U} \cap \mathrm{C} ^ {\prime}) = \mathrm{U} \cap w ^ {\prime} (\mathrm{C} ^ {\prime});
\]

since the union of the \(w'(C)\) for \(w' \in W_a\) is dense in \(E\) , the union of the \(U \cap w'(C) = U \cap w'(C')\) is dense in \(U\) , and the union of the \(w'(\overline{C})\) for \(w' \in W_a\) thus contains \(U\) . Finally, if \(H\) is a wall of \(C'\) , there exist a point \(c \in U \cap H\) and an open neighbourhood \(V \subset U\) of \(c\) such that \(V \cap C'\) is the intersection of \(V\) and the open half-space bounded by \(H\) containing \(C'\) ; since \(V \cap C' = V \cap U \cap C' = V \cap U \cap C = V \cap C\) , it follows that \(H\) is a wall of \(C\) . If \(W\) is infinite and irreducible, \(C\) is an open simplex (no. 9, Prop. 8) and so has one more wall than the open simplicial cone \(C'\) .

COROLLARY. Assume that W is essential.

\begin{enumerate}
\def\labelenumi{(\roman{enumi})}
\item
  If \(a \in \mathbb{E}\) is special, there exists a chamber \(\mathbb{C}\) such that \(a\) is an extremal point of \(\overline{\mathbb{C}}\) .
\item
  If \(\mathrm{C}\) is a chamber, there exists an extremal point of \(\overline{\mathrm{C}}\) that is special.
\end{enumerate}

The first assertion follows from Prop. 11. The second follows from the first and the fact that W acts transitively on the set of chambers.

However, an extremal point of \(\overline{C}\) is not necessarily a special point for W (cf.~Chap. VI, Plate X, Systems \(B_{2}\) and \(G_{2}\) ).

Remark. 1) Assume that W is essential, irreducible and infinite and retain the notation of Prop. 11. Since U is an isomorphism from \(W_{a}\) to \(U(W)\) , it follows that the Coxeter graph of the group of displacements \(U(W)\) (which is generated by the reflections \(U(s_{H})\) for \(H \in \mathfrak{H}\) ) can be obtained from the Coxeter graph of W by omitting the vertex i corresponding to the unique wall of C that is not a wall of \(C'\) .

PROPOSITION 12. Assume that W is essential. Let a be a special point, let L(a) be the set of its transforms under the group of translations L, and let C be a chamber. Then \(\overline{C}\) meets L(a) in a unique point; this point is an extremal point of \(\overline{C}\) .

There exists a chamber \(C_{1}\) such that a is an extremal point of \(\overline{C}_{1}\) (Cor. of Prop. 11). Every chamber is of the form \(\mathbf{C}=tw^{\prime}(\mathbf{C}_{1})\) with \(w^{\prime}\in W_{a}\) and \(t\in L\) since \(W=W_{a}.L\) ; thus \(tw^{\prime}(a)=t(a)\in\mathrm{L}(a)\) is an extremal point of \(\overline{C}\) . On the other hand, \(\overline{C}\) cannot contain two distinct points of \(\mathrm{L}(a)\) since \(\overline{C}\) is a fundamental domain for W (no. 3, Th. 2).

Remark. 2) The set L(a) is contained in the set of special points, but in general is distinct from it (cf.~Chap. VI, §2, no. 2 and Plates I to VI).

